% Page 011

\seite{11}

bis Schönecken, Kreis Prüm, gekommen sein.\ob
Die Kisten kamen mit Quarz angefüllt in\ob
Sefferweich an und die Verwandten hatten das\ob
Nachsehen.

\margmark{Gehört zu\\Ritter von der Els}%
   Daß die oben angeführten, der Pfarrkirche zu Seffern geschenkten\ob
Güter des von der Els nicht unbedeutend waren, beweist\ob
folgende Thatsache: Ein Theilungsvertrag aus dem Jahre\ob
1610, aufgenommen zu Wavern vor Schultheiß und\ob
Schöffen des Hofes Seffern. In diesem Theilungsver\ohb
trage setzten sich die beiden Häuser \enquote{Trauten} u\ob
\enquote{Roths}, welche bis zum Jahre 1610 nur einen Stock\ob
bildeten, derart auseinander, daß Roths die Besitzungen\ob
diesseits der Dorfstraße und Trauten jene jenseits\ob
derselben erhalten sollten. Die Dorfstraße bildete so\ohb
mit die Grenze zwischen beiden Theilen. Nun muß\ob
man aber wissen, daß die Dorfstraße früher eine\ob
andere Lage hatte, als die heutige. Die alte Dorfstraße\ob
ging früher am Brunnen und an der Kapelle vorbei.\ob
Die jetzige Dorfstraße wurde erst im Jahre 1848 oder\ob
1849 durch Betreiben des damaligen Bürgermeisters\ob
\margmark{Ehlenz\\größtenteils\\Dem\\Pastor\\Seffern von\\Seffern!}%
Axer\unsicher, eines geborenen Bitburgers, hergestellt, dem\ob
es auch \sout{[...]} größtenteils zu verdanken ist, daß im Jahre\ob
1851 der neue Kirchenbau in Seffern in Angriff\ob
genommen wurde. Alle Anzeigen deuten darauf\ob
hin, daß anfänglich \enquote{Trautenhaus} bis zum Jahre 1810\ob
und von da ab die beiden Häuser Trauten u Roths die
\pend
